\subsection{开源与闭源的分岔：一个正在发生的商业对照}

本文论证至此的理论框架，在当前的AI产业格局中已经可以观察到具体的对照案例。

\subsubsection{模式一：开源直通车——GLM与DeepSeek}

中国的开源大模型生态在过去两年中迅速发展。以智谱AI的GLM系列和深度求索的DeepSeek系列为代表的开源模型，允许任何企业直接下载模型权重、在自有服务器上完成本地部署、根据特定业务需求进行微调。这一路径的核心特征是\textbf{去中介化}：无需第三方审批、数据主权归属企业自身、无按调用量计费的边际成本、企业可针对自身业务场景自由微调。

这种模式在制造业、农业、县域政务等非数字原生领域具有特殊的赋能价值。一家宁阳的农产品加工企业，不需要通过硅谷私募基金的AI审计，就可以部署开源模型来优化供应链管理。这正是本文所论证的光谱爬升入口对所有网络连接者开放在产业层面的具体体现。

\subsubsection{模式二：闭源把关——OpenAI与私募基金的AI审计}

与之形成对照的是闭源生态中正在出现的一种新型商业模式。据报道和分析，OpenAI与部分私募股权基金建立了合作关系，由基金对其投资组合中的企业进行AI转型就绪度评估，评估通过的企业获得OpenAI工具的优先部署权和基金追加投资。这一模式的逻辑可以概括为\textbf{再中介化}：AI能力不是企业可以自主获取的生产要素，而是需要经过外部机构的审核和授权；AI工具的获取与融资资格捆绑；企业对自身AI就绪度的判断权被让渡给外部评估者；使用闭源API意味着企业数据需经过第三方服务器。

这一模式在金融逻辑上是合理的——私募基金需要评估被投企业的技术能力以管理风险。但从认知民主化的角度看，它恰恰构成了本文所警示的认知分层固化机制：AI能力不是普惠的生产要素，而是需要许可才能获取的特权。

\subsubsection{对照的启示}

两种模式的并存使得本文的理论论证具有了现实紧迫性。开源路径证明了去中介化的AI赋能是可行的——企业可以自主获取、部署和定制AI工具，无需通过守门人。闭源路径则展示了如果没有开源替代方案，AI能力将如何被金融中介重新围墙化。这一对照也是对市场会解决可及性问题的经验回应：市场确实使闭源AI的价格下降，但价格下降不等于去中介化。一个便宜的、但需要通过私募基金审计才能使用的AI，仍然是守门人控制下的AI。认知民主化需要的不是更便宜的许可，而是不需要许可的能力。
